\section{The finite scale induction}
\label{sec:bootstrap}

The preceding depth gives localization for projections with small
support. We now extend that estimate to arbitrary product projections.
The extension uses the inverse circuit at a larger mismatch scale.
The main constraint is how far that scale can move: to prove
$\cB_d(s,t)$ from $\cB_d(t,b)$, we need
$t<b<2t-s$. This interval leaves room for a polynomial whose
approximation error decays faster than $e^{-N^s}$ while its degree
and the logarithm of its weighted coefficient norm remain below the
mismatch threshold $N^t$.

\begin{lemma}[One backward step]
\label{lem:backward}
Fix \(d\ge1\), and assume \(\cB_{d-1}(s,t)\) for every fixed
\(0\le s<t\). Let \(s,t,b\) be fixed real numbers with
\[
 0\le s<t<b<2t-s.
\]
If \(\cB_d(t,b)\) holds, then \(\cB_d(s,t)\) holds.
\end{lemma}

\begin{proof}
Choose fixed auxiliary exponents
\[
 s<a<2t-b,\qquad \frac{a+b}{2}<u<w<t.
\]
These choices are possible because $b<2t-s$. Since $b>t>a$,
they give exactly the inequalities required by the polynomial and
small-support estimates:
\begin{equation}
\label{eq:step-parameters}
 0\le s<a<u<w<t<b,\qquad u>\frac{a+b}{2}.
\end{equation}
Fix \(U\in\cU_d(N)\) and counts \(M,D\), and put
\[
 P=[M=0],\qquad H=[M\ge N^t],\qquad
 C=UDU^\dagger,\qquad Q=[C=0]=U[D=0]U^\dagger.
\]
Let \(p(x)=\sum_kc_kx^k\) be the polynomial in
\Cref{lem:polynomial} for \(a,b,u\). We will first control
$Hp(C)P$ by small-support localization. Squaring the desired norm
will then place the remaining high-$C$ error on $PQ$, where the
larger-scale hypothesis applies after inversion.

Each count summand is a rank-one projection on one qubit.
Expanding \(D^k\) gives at most \(N^k\) ordered monomials, each a
product projection on at most \(k\) qubits: repeated indices collapse
by idempotence, while different indices commute.
For \(k\le\deg p\le N^u\), \Cref{lem:small-support} applies to
every such monomial with the same fixed exponents \(u,w,t\).
Since \(p(C)=Up(D)U^\dagger\), its coefficient bound gives
\begin{equation}
\label{eq:polynomial-transition}
 \norm{Hp(C)P}
 \le e^{-N^w}\sum_k|c_k|N^k
 \le e^{N^u-N^w}.
\end{equation}
The constant term causes no difficulty: it is a multiple of \(I\),
and \(HP=0\). Also, \(\norm C\le N\) gives
\begin{equation}
\label{eq:polynomial-norm}
 \norm{p(C)}\le\sum_k|c_k|N^k\le e^{N^u}.
\end{equation}

The polynomial need not approximate \(Q\) on all of \(\mathcal H_N\).
Squaring the desired norm makes the approximation error occur as
\((Q-p(C))PQ\), rather than \((Q-p(C))P\):
\begin{align}
\norm{HQP}^2
 &=\norm{HQPQH}\notag\\
 &\le \norm{Hp(C)PQH}+\norm{H(Q-p(C))PQH}\notag\\
 &\le \norm{Hp(C)P}+\norm{(Q-p(C))PQ}.
\label{eq:square}
\end{align}
Only the leftmost \(Q\) was replaced in this decomposition. No commutation
of \(P\) with \(C\) or \(Q\) is involved.

Split the last error at the spectral threshold \(N^b\) of \(C\).
The spectral projections of \(C\) commute with \(Q\) and \(p(C)\).
On \([C<N^b]\), the error is zero at eigenvalue \(0\), since
\(p(0)=1\), and at most \(e^{-N^a}\) at every positive eigenvalue,
since these eigenvalues are integers. Hence
\begin{equation}
\label{eq:error-split}
 \norm{(Q-p(C))PQ}
 \le e^{-N^a}+(1+e^{N^u})\norm{[C\ge N^b]PQ}.
\end{equation}
The remaining factor is exactly the transition controlled by the
inverse circuit. Unitary conjugation gives
\begin{align}
 \norm{[C\ge N^b]PQ}
 &=\norm{[D\ge N^b]\,U^\dagger[M=0]U[D=0]}\notag\\
 &\le e^{-N^t}.
\label{eq:inverse-tail}
\end{align}
Here \(\cB_d(t,b)\) applies with circuit \(U^\dagger\), mismatch
count \(D\), and projected count \(M\). Inverse closure
\eqref{eq:inverse} and the universal quantifiers in
\(\cB_d(t,b)\) are both used.

Combining \eqref{eq:polynomial-transition}, \eqref{eq:square},
\eqref{eq:error-split}, and \eqref{eq:inverse-tail} yields
\begin{equation}
\label{eq:step-bound}
 \norm{HQP}^2
 \le e^{N^u-N^w}+e^{-N^a}+(1+e^{N^u})e^{-N^t}.
\end{equation}
By \eqref{eq:step-parameters}, each term is eventually at most
\(\tfrac13e^{-2N^s}\). Indeed \(w>u\), \(t>u\), and
\(a,w,t>s\). This remains true for \(s=0\), when the comparison
quantity is a fixed positive constant. Taking the square root proves
\(\cB_d(s,t)\). All sufficient lower bounds on \(N\) depend only
on the fixed parameters and \(d\), so the conclusion is uniform.
\end{proof}

We now choose a finite chain of scales within this interval. Each step
will advance the mismatch exponent by almost the current gap $t-s$.
Keeping at least half of the initial gap for $O(1/(t-s))$ steps will
carry that exponent above one, where the tail projection vanishes.
A slack of order $(t-s)^2$ per step is small enough to preserve the gap.

\begin{proof}[Proof of \Cref{thm:localization}]
The case \(d=0\) is \Cref{lem:base}. Fix \(d\ge1\), and assume the
theorem at depth \(d-1\) for every fixed exponent pair.
It remains to prove \(\cB_d(s,t)\) for fixed \(0\le s<t\le1\);
the case \(t>1\) is immediate.
Put
\[
 g=t-s,\qquad \eta=\frac{g^2}{100},\qquad
 s_0=s,\quad t_0=t,
\]
and define
\begin{equation}
\label{eq:scale-recurrence}
 s_{i+1}=t_i,\qquad t_{i+1}=2t_i-s_i-3\eta.
\end{equation}
Writing \(g_i=t_i-s_i\), we have
\[
 g_i=g-3i\eta,\qquad t_{i+1}-t_i=g_{i+1}.
\]
Let \(K=\lceil2/g\rceil\). Since \(0<g\le1\),
\[
 3K\eta\le \frac{3(2/g+1)g^2}{100}
 \le0.09g<\frac g2.
\]
Thus \(g_i>g/2\) for \(i\le K\), and
\[
 t_K=t+\sum_{i=1}^K g_i>t+Kg/2\ge t+1>1.
\]
Let \(J\le K\) be the first index for which \(t_J>1\).
This index and every exponent in the chain are independent of \(N\).
The assertion \(\cB_d(s_J,t_J)\) is trivial.

\begin{figure}[htbp]
\centering
\[
\underbrace{\cB_d(s_0,t_0)}_{\text{target}}
\ \Longleftarrow\
\cB_d(s_1,t_1)
\ \Longleftarrow\ \cdots\ \Longleftarrow\
\underbrace{\cB_d(s_J,t_J)}_{t_J>1:\ \text{trivial}} .
\]
\caption{The same-depth induction runs from right to left along a
finite chain; \(J\) is the first index with \(t_J>1\).
Each implication is \Cref{lem:backward} and also uses
the already established localization theorem at depth \(d-1\).
The inverse circuit is substituted only into the universal statement
at the next pair, after that statement has been proved.}
\label{fig:scales}
\end{figure}

Suppose now that $\cB_d(s_{i+1},t_{i+1})$ has been proved,
where $i<J$. The recurrence and the positive next gap give
\[
 t_i<t_{i+1}=2t_i-s_i-3\eta<2t_i-s_i.
\]
Thus \Cref{lem:backward} applies with $b=t_{i+1}$. Its assumed
assertion is precisely
$\cB_d(t_i,t_{i+1})=\cB_d(s_{i+1},t_{i+1})$, so it proves
$\cB_d(s_i,t_i)$.
Working backward reaches the target pair, as in \Cref{fig:scales}.

There are only finitely many steps for each fixed \(d,s,t\).
Taking the maximum of their sufficient lower bounds on \(N\)
preserves uniformity over every circuit and pair of counts. Each step
uses only finitely many fixed exponent tuples from depth $d-1$; the
number of monomials and their varying supports are paid for by the
explicit coefficient estimates, not by further asymptotic thresholds.
Thus the induction does not require a threshold uniform over all real
exponent pairs. This completes the depth induction.
\end{proof}
